\section{Discussion}

% OWNER: Sanjana Garimella (discussion per proposal)

\begin{frame}{When does the rule break?}
  \begin{itemize}
    \item \textbf{Threshold, not gradual}: absorbs noise up to $10\%$, snaps at $50\%$
      --- like a child's grammar surviving occasional speech errors.
    \item \textbf{Attractor gap is the early warning}: it climbs ($0.07 \to 0.20$)
      while overall accuracy still looks healthy.
    \item \textbf{Delay vs.\ prevention}: low noise only \emph{delays} acquisition
      (curves converge); $50\%$ noise \emph{prevents} it and entrenches the linear rule.
  \end{itemize}
\end{frame}

\begin{frame}{Does structure buy robustness?}
  \begin{itemize}
    \item The headline question --- does a syntactic prior (RNNG) resist noise
      better than a sequential LSTM? --- needs the RNNG runs (in progress).
    \item \textbf{LSTM baseline}: no explicit hierarchy bias $\Rightarrow$ learns the
      rule cleanly when input is clean, but drifts to the linear shortcut as noise grows.
    \item If RNNG holds the hierarchical rule where the LSTM has caved
      $\Rightarrow$ structure buys robustness.
    \item If both fail together $\Rightarrow$ the bottleneck is input quality, not
      architecture --- speaking to poverty-of-the-stimulus.
  \end{itemize}
\end{frame}

\begin{frame}{Takeaways and future work}
  \textbf{Takeaways}
  \begin{itemize}
    \item LSTMs tolerate moderate SVA noise but have a sharp breaking point.
    \item Noise trades the hierarchical rule for the linear nearest-noun heuristic.
    \item Heavy noise prevents acquisition outright --- more training makes it worse.
  \end{itemize}
  \textbf{Limitations \& next steps}
  \begin{itemize}
    \item Synthetic, uniform noise; templatic corpus; single-seed trajectories.
    \item Complete the RNNG arm; try ON-LSTM as a no-gold-tree syntactic baseline.
    \item Move toward naturalistic, frequency-dependent noise and larger scale.
  \end{itemize}
\end{frame}
